\documentclass[12pt,a4paper]{article}

% Pacotes essenciais
\usepackage[utf8]{inputenc}
\usepackage[T1]{fontenc}
\usepackage[brazilian]{babel}
% Removido pacote graphicx que causava erro e não é mais necessário
\usepackage{indentfirst} % Corrigido de "indentof" para "indentfirst"
\usepackage{setspace}
\usepackage[left=3cm,right=2cm,top=3cm,bottom=2cm]{geometry}
\usepackage{hyperref}
\usepackage{cite}


\onehalfspacing
\setlength{\parindent}{1.25cm}

\begin{document}


\begin{titlepage}
    \centering
    \vspace*{1cm}
    
    \textbf{\Large UNIVERSIDADE FEDERAL DO SEMI-ÁRIDO (UFERSA)}\\[0.3cm]
    \textbf{\large CURSO: INTERDICIPLINAR EM TECNOLOGIA DA INFORMAÇÃO}\\[0.3cm]
    \textbf{\large DISCIPLINA: TESTE DE SOFTWARE}\\[2cm]
    
    \textbf{ABNOAN GABRIEL FERREIRA DA SILVA}\\
    \textbf{JÉSSICA ISABELA CARDOSO DE CASTRO}\\
    \textbf{JOSÉ DANILO SANTOS DO NASCIMENTO}\\[4cm]
    
    \textbf{\Large EconoQuiz: Uma Plataforma Gamificada para a Promoção da Literacia Financeira e o Suporte à ODS 8 (Trabalho Decente e Crescimento Econômico)}\\[4cm]
    
    \vfill
    
    \textbf{PAU DOS FERROS - RN}\\
    \textbf{2025}
\end{titlepage}


\section*{RESUMO}
\addcontentsline{toc}{section}{RESUMO}

Este artigo apresenta o desenvolvimento e a arquitetura do EconoQuiz, uma aplicação de quiz gamificada projetada para aumentar a literacia financeira de forma acessível e envolvente. O projeto se alinha à Meta 8.10 do Objetivo de Desenvolvimento Sustentável (ODS) 8: Trabalho Decente e Crescimento Econômico, buscando fortalecer a inclusão financeira e o acesso a serviços econômicos básicos. O desempenho dos usuários é avaliado de maneira dinâmica, utilizando um sistema de pontuação que reflete o nível de proficiência em tópicos financeiros essenciais, com foco em jovens e adultos. Detalham-se os requisitos funcionais e não funcionais que garantem uma experiência de aprendizado robusta, com partidas baseadas em níveis de dificuldade e pontuações atualizadas em tempo real em um ranking competitivo. A arquitetura da plataforma segue o padrão Model-View-Controller (MVC), implementado com React Native no front-end e Node.js/Express no back-end. O EconoQuiz é proposto como uma solução tecnológica escalável para o fortalecimento da educação financeira e o desenvolvimento sustentável.

\noindent\textbf{Palavras-chave:} Literacia Financeira; Gamificação; ODS 8; Desenvolvimento Sustentável; Tecnologia Educacional.



% Conteúdo
\section{INTRODUÇÃO}

A literacia financeira é um pilar essencial para a autonomia e o empoderamento econômico na sociedade contemporânea. Segundo o Core Competencies Framework da OCDE/INFE, a literacia financeira envolve competências para compreender conceitos e riscos e aplicar esse conhecimento em decisões cotidianas (OECD INFE, 2018). No Brasil, dados e materiais do Banco Central mostram lacunas expressivas em conhecimentos sobre juros, inflação e orçamento, especialmente entre jovens (BANCO CENTRAL DO BRASIL, 2018). Em resposta a essa lacuna educacional, o presente trabalho descreve a concepção e a arquitetura do EconoQuiz, uma aplicação móvel e web desenhada para transformar a educação financeira em uma experiência envolvente e acessível.

O EconoQuiz surge como uma ferramenta de apoio à inclusão financeira, alinhando-se diretamente à Meta 8.10 do Objetivo de Desenvolvimento Sustentável (ODS) 8 da ONU, que foca na promoção do trabalho decente e do crescimento econômico (IPEA, 2025). Por meio da capacitação em habilidades financeiras, o projeto visa aprimorar a capacidade de jovens e adultos para uma participação plena e informada no mercado.

O sistema opera com base em um mecanismo de quiz gamificado. As partidas variam em dificuldade e atribuem pontuações que definem o nível de proficiência e a posição do usuário em um ranking global. O design é suportado por um conjunto rigoroso de requisitos funcionais e não funcionais para garantir escalabilidade e alto desempenho. Este artigo detalhará a fundamentação teórica que justifica a metodologia de gamificação e apresentará o design técnico da solução como um agente de mudança social e econômica.

\section{DESENVOLVIMENTO}

\subsection{Fundamentação Teórica}

O design do EconoQuiz baseia-se na interseção entre dois eixos principais: a literacia financeira como meio de inclusão social e o uso da gamificação como estratégia pedagógica para promover o engajamento e a retenção de conhecimento.

\subsubsection{Literacia Financeira como Fator de Inclusão}

A literacia financeira é um componente central da inclusão econômica e social. Pesquisas apontam que indivíduos com maiores habilidades financeiras tomam decisões mais conscientes, evitam superendividamento e planejam melhor sua vida econômica (LUSARDI; MITCHELL, 2014).

O EconoQuiz busca responder a essa demanda por meio da disseminação de conceitos básicos de economia, orçamento pessoal e consumo consciente. O jogo utiliza uma abordagem progressiva, em que as perguntas são classificadas em três níveis de dificuldade: fácil, médio e difícil, representando a evolução cognitiva do usuário. Dessa forma, o aplicativo não apenas transmite conhecimento, mas também mensura o progresso e estimula o aprendizado contínuo.

Ao alinhar-se à ODS 8, o projeto contribui para o fortalecimento da capacidade financeira individual, elemento indispensável para a promoção do trabalho decente e o crescimento econômico sustentável. A inclusão financeira é, portanto, não apenas uma meta social, mas também um instrumento de equidade e desenvolvimento (IPEA, 2025).

\subsubsection{A Gamificação como Estratégia de Aprendizagem}

A gamificação consiste na utilização de elementos de jogos, como desafios, recompensas, níveis e rankings, em contextos não relacionados a jogos, visando aumentar a motivação e o engajamento dos usuários e tem suporte empírico consistente para aumentar engajamento e retenção de conteúdo (HAMARI; KOIVISTO; SARSA, 2014). Revisões sistemáticas mostram ganhos em motivação, persistência e aprendizagem significativa (DOMÍNGUEZ et al., 2013).

No contexto do EconoQuiz, a gamificação é aplicada através de:

\begin{itemize}
    \item \textbf{Sistema de Pontuação:} baseado na dificuldade da questão e na precisão das respostas, refletindo a progressão de conhecimento.
    \item \textbf{Ranking Global:} incentivo à competição saudável e à socialização entre usuários.
    \item \textbf{Feedback Imediato:} reforço positivo que orienta o aprendizado e estimula a correção de erros.
    \item \textbf{Dicas Contextuais:} auxiliam na compreensão de temas complexos, reforçando o caráter educativo do jogo.
\end{itemize}

Essa estrutura permite que o EconoQuiz vá além do entretenimento, posicionando-se como uma ferramenta pedagógica eficaz que alia teoria econômica e prática interativa.

\section{DESENVOLVIMENTO DA PLATAFORMA}

A metodologia de desenvolvimento do EconoQuiz seguiu uma abordagem ágil, garantindo que o produto final atendesse aos requisitos de usabilidade e performance. A arquitetura foi concebida para suportar um grande volume de acessos e garantir a integridade dos dados, adotando um modelo de três camadas.

\subsection{Requisitos de Sistema}

Os requisitos guiaram o desenvolvimento e foram classificados em:

\begin{itemize}
    \item \textbf{Funcionais (RF):} Focados na interação do usuário, como o sistema robusto de login/cadastro, a geração de quizzes parametrizados por nível de dificuldade, o algoritmo de pontuação em tempo real e a funcionalidade de gestão e exibição do ranking global.
    \item \textbf{Não Funcionais (RNF):} Relacionados à qualidade da plataforma, enfatizando a Usabilidade (interface clara e responsiva), a Segurança dos dados pessoais e financeiros, e a Performance (baixa latência para garantir a fluidez das partidas).
\end{itemize}

\subsection{Design Tecnológico e Padrão MVC}

A plataforma adota o padrão Model-View-Controller (MVC), que assegura a separação de responsabilidades entre as camadas do sistema, facilitando a manutenção e a escalabilidade (GIGASYSTEMS, 2015).

\begin{itemize}
    \item \textbf{Front-end (View):} Utiliza React Native para o desenvolvimento do aplicativo móvel, permitindo a compatibilidade nativa com iOS e Android a partir de um único codebase (REACT NATIVE, 2025). O design de interface incorpora elementos de jogos para reforçar o engajamento.
    \item \textbf{Back-end (Controller e Model):} É construído com Node.js e Express.js. Esta escolha garante um processamento assíncrono eficiente, vital para gerenciar as sessões de quiz e as constantes atualizações do ranking. O Banco de Dados (BD) é central para o armazenamento de perfis de usuário, o acervo de questões categorizadas e o registro histórico das pontuações.
\end{itemize}

\section{RESULTADOS E DISCUSSÃO}

Embora este artigo se concentre na concepção e arquitetura do EconoQuiz, os resultados esperados e a discussão são fundamentados nas melhores práticas da literatura de Tecnologia Educacional e Gamificação. O desempenho da plataforma é analisado sob três eixos: eficácia pedagógica, desempenho técnico e impacto social.

\subsection{Eficácia da Metodologia Gamificada}

Espera-se que a mecânica de pontos, níveis e ranking do EconoQuiz proporcione um aumento significativo no engajamento dos usuários. Estudos mostram que metodologias gamificadas elevam significativamente o engajamento do estudante (HAMARI; KOIVISTO; SARSA, 2014). A estrutura do EconoQuiz, com metas claras, recompensas e feedback, favorece retenção de conteúdo e aprendizagem ativa (DOMÍNGUEZ et al., 2013).

O sistema de pontuação funciona como uma métrica de proficiência, permitindo aos usuários não apenas avaliar seu conhecimento, mas também traçar um caminho claro para o desenvolvimento. A discussão central é que a diversão fornecida pelo design do jogo atua como um potente catalisador para a aprendizagem autônoma e continuada em um tema de alta relevância social.

\subsection{Desempenho e Escalabilidade Técnica}

O uso da stack React Native/Node.js/Express.js, aliada ao padrão MVC, projeta o EconoQuiz para ser altamente escalável. O modelo assíncrono do Node.js permite que o servidor gerencie milhares de conexões simultâneas (partidas de quiz) com baixa latência, o que é crucial para uma experiência de usuário fluida e responsiva, especialmente na atualização em tempo real do ranking. O desempenho do sistema é, portanto, diretamente coerente com o objetivo de alcançar a Meta 8.10 do ODS 8, que exige uma solução tecnológica capaz de atender a uma ampla base de usuários para promover a inclusão financeira massiva.

\subsection{Impacto Social e ODS 8}

O impacto do EconoQuiz é medido pela sua contribuição para a Meta 8.10. Ao tornar o aprendizado financeiro acessível em dispositivos móveis, a plataforma atinge populações que historicamente enfrentam barreiras ao acesso à informação econômica (IPEA, 2025). A discussão é que a capacitação gerada pelo EconoQuiz transcende a aquisição de conhecimento, promovendo a capacidade de tomar decisões financeiras conscientes, o que é um pré-requisito fundamental para a participação plena no mercado de trabalho e para a redução da vulnerabilidade econômica, fortalecendo a resiliência financeira.



% Referências
\begin{thebibliography}{99}

\bibitem{bc2018}
BANCO CENTRAL DO BRASIL. \textbf{Caderno de Cidadania Financeira — Cuidando do seu dinheiro: Gestão de finanças pessoais.} Brasília: Banco Central do Brasil, 2018. Disponível em: \url{https://www.bcb.gov.br/content/cidadaniafinanceira/documentos_cidadania/Cuidando_do_seu_dinheiro_Gestao_de_Financas_Pessoais/caderno_cidadania_financeira.pdf}. Acesso em: 9 nov. 2025.

\bibitem{dominguez2013}
DOMÍNGUEZ, A. et al. Gamifying learning experiences: practical implications and outcomes. \textbf{Computers \& Education}, v. 63, p. 380–392, 2013. Disponível em: \url{https://doi.org/10.1016/j.compedu.2012.12.020}. Acesso em: 14 nov. 2025.

\bibitem{hamari2014}
HAMARI, J.; KOIVISTO, J.; SARSA, H. Does Gamification Work? — A Literature Review of Empirical Studies on Gamification. \textbf{Proceedings of the 47th Hawaii International Conference on System Sciences (HICSS)}, 2014. Disponível em: \url{https://creativegames.org.uk/modules/Gamification/Hamari_etal_Does_gamification_work-2014.pdf}. Acesso em: 15 nov. 2025.

\bibitem{gigasystems2015}
GIGASYSTEMS. \textbf{Arquitetura MVC de forma simples.} Gigasystems, 07 jan. 2015. Disponível em: \url{https://www.gigasystems.com.br/artigo/59/arquitetura-mvc-de-forma-simples}. Acesso em: 15 nov. 2025.

\bibitem{ipea2025}
INSTITUTO DE PESQUISA ECONÔMICA APLICADA (IPEA). \textbf{ODS 8 – Trabalho Decente e Crescimento Econômico.} Brasília: Ipea, [s.d.]. Disponível em: \url{https://www.ipea.gov.br/ods/ods8.html}. Acesso em: 9 nov. 2025.

\bibitem{lusardi2014}
LUSARDI, A.; MITCHELL, O. S. The Economic Importance of Financial Literacy: Theory and Evidence. \textbf{Journal of Economic Literature}, v. 52, n. 1, p. 5–44, 2014. Disponível em: \url{https://gflec.org/wp-content/uploads/2014/12/economic-importance-financial-literacy-theory-evidence.pdf}. Acesso em: 15 nov. 2025.

\bibitem{oecd2018}
OECD INFE. \textbf{OECD/INFE Core Competencies Framework on Financial Literacy.} Paris: OECD, 2018. Disponível em: \url{https://www.oecd.org/finance/financial-education/core-competencies.htm}. Acesso em: 9 nov. 2025.

\bibitem{reactnative2025}
REACT NATIVE. \textbf{React Native — Getting started / Documentation.} Meta. Disponível em: \url{https://reactnative.dev/docs/getting-started}. Acesso em: 9 nov. 2025.

\end{thebibliography}

\end{document}
